\section{Background}
\label{sec:background}

Three concepts are load-bearing for the method sections that follow.

\textbf{Variant allele frequency (VAF).}
The fraction of sequencing reads at a locus that carry the alternate allele.
In a tumour sample, VAF depends on tumour purity, ploidy, cancer cell fraction, and copy number.
A heterozygous SNV in a pure diploid tumour has an expected VAF of 0.50.
At 60\% purity, the same variant has an expected VAF of 0.30.

\textbf{Unique molecular identifiers (UMIs).}
Short random sequences (typically 4 to 14 bp) ligated to each DNA molecule before PCR amplification.
After sequencing, reads sharing the same UMI and alignment position are grouped into families originating from one template molecule.
Consensus calling across a family suppresses PCR and sequencing errors.
Duplex sequencing extends this by tagging both strands of a double-stranded molecule independently, requiring agreement between strand-specific consensus reads before reporting a variant.

\textbf{Cell-free DNA (cfDNA).}
Short DNA fragments circulating in blood plasma, released by apoptosis and necrosis.
Normal cfDNA has a characteristic fragment length distribution dominated by a mono-nucleosomal peak at ${\sim}167$\,bp (one nucleosome plus linker), with sub-nucleosomal and di-nucleosomal components.
Tumour-derived cfDNA (ctDNA) fragments are shorter on average, peaking at ${\sim}143$--$147$\,bp.
This size difference is exploited by fragment-based liquid biopsy tools.
